\code{LWPTracking} is used to configure the handling of LWP events for its
associated Process.

\begin{apient}
static void setDefaultTrackLWPs(bool b)
\end{apient}

\apidesc{
Sets the default handling mechanism for LWP events across all Process objects to
interrupt-driven (b = true) or polling-driven (b = false).
}

\begin{apient}
static bool getDefaultTrackLWPs()
\end{apient}

\apidesc{
Returns the current default handling mechanism for LWP events across all Process
objects. A return value of true indicates interrupt-driven while false
indicates polling-driven.
}

\begin{apient}
bool setTrackLWPs(bool b) const
\end{apient}

\apidesc{
Sets the LWP event handling mechanism for the associated Process object to
interrupt-driven (b = true) or polling-driven (b = false).
Returns true on success and false on failure.
}

\begin{apient}
bool getTrackLWPs() const
\end{apient}

\apidesc{
Returns the current LWP event handling mechanism for the associated Process
object. A return value of true indicates interrupt-driven while false
indicates polling-driven.
}

\begin{apient}
bool refreshLWPs()
\end{apient}

\apidesc{
Manually polls for queued LWP events to handle.
Returns true on success and false on failure.
}
